\subsection{Deterministic histories must change the fitting population}
A simulator can enforce a structural rule while estimating the stochastic
part of that rule on the wrong records. The selected pygformula documentation
describes a covariate restriction as a known value outside specified histories.
It directs the fitting step to exclude those deterministic histories, then
directs simulation to assign their known value. For a bank, eligibility may
make a service-completion indicator identically zero: the completion model
then describes eligible accounts, rather than an average that includes
accounts unable to receive the service. This is a distinction between two
conditional distributions, even before any causal effect is considered.
\citep{rel25pygfit}

The documented input is a list of restrictions, each beginning with the
covariate's name. In the selected fitting helper, the list used to identify
restricted covariates repeats the whole first restriction instead of extracting
each name. Its subsequent membership test compares a name with those lists,
so the documented input does not enter the intended filtering branch.
The simulation helper does extract the names and assigns deterministic
values on the corresponding histories. The constructor forwards the supplied
restrictions directly. Thus the two steps implement different populations
on this selected path; this conclusion follows from inspected source, not
from an executed package experiment. It is specific to this revision and
schema, and does not assert that every fitted model changes numerically.
\citep{rel25pygfit}

An exact constructed likelihood shows when the distinction matters. Let
$E$ indicate eligibility, $Z$ completion and $X$ a binary predictor. At a
post-baseline fitting time, each value of $X$ has ten eligible records,
eight with $Z=1$ and two with $Z=0$, and ten ineligible records with the known
value $Z=0$. Use the logistic model
\[
 \operatorname{logit}\Pr(Z=1\mid X,E=1)=\beta_0+\beta_1X.
\]
Restricting estimation to eligible records gives fitted probability $4/5$
in both predictor cells. Pooling the deterministic zeros gives $2/5$ in both.
For either fit, the intercept and slope scores vanish at the corresponding
constant probability; both cells contain successes and failures, and the
information matrix is positive definite. These are finite, unique likelihood
solutions, not an appeal to an intercept-only formula that the input parser
might handle differently. The predictor is deliberately uninformative within
eligibility; no claim is made about a model that separately represents every
eligibility history.~\citep{rel25pygwitness}

If half the simulated accounts are eligible and the other half are correctly
set to zero, the completion mean is $2/5$ under restricted estimation but
$1/5$ under the pooled fit. A stipulated USD100 receipt contingent on completion
and USD 30 cost per account give expected margins of USD 10 and USD $-10$.
This is a contract-payoff calculation under an invented law, not an estimated
causal service benefit. Increasing simulation size converges more precisely
to the respective fitted laws; it cannot repair their different fitting
populations. The proposed worker should therefore retain the exact fitting
row set for each transition, verify deterministic assignments separately, and
withhold this restriction profile until a corrected native path reproduces
the declared law. An explicitly fitted stochastic-subpopulation model is a
substantive alternative, provided that its conditioning histories and
standardization population remain visible.~\citep{rel25pygwitness,rel25pygfit}

\subsection{Censoring labels and competing risks define different denominators}
The input clock places covariates on interval $k$ and the outcome on
interval $k+1$. The selected outcome helper uses nonmissing outcomes at
nonnegative times; with explicit competing events it also requires no current
competing event. A censoring-name argument is not passed to that helper.
The censoring model instead supports the observed natural-course inverse
probability weighted comparator. Supplying that model therefore does not
automatically select the outcome likelihood rows. The documentation permits
missing or zero outcomes on a survival censoring row under its stated
fine-interval qualification; for an end-of-follow-up outcome it requests a
missing value. These conventions concern which event information is observed,
not merely how a missing field is displayed.~\citep{rel25pygfit}

For a deliberately coarse interval, suppose 100 accounts supply eight observed
defaults, 72 observed nondefaults and 20 outcomes unknown after censoring.
A constant outcome likelihood over observed outcomes gives $8/80=1/10$;
coding the unknown outcomes zero gives $8/100=2/25$. A stipulated contract
paying USD 10,000 on default and charging USD 900 has model-based margins
USD 100 and USD $-100$, respectively. Neither choice alone identifies the
complete finite-population risk: the 20 unknown outcomes permit a proportion
between $2/25$ and $7/25$. Interpreting the observed-outcome fit as a target
risk additionally requires an appropriate conditional independent-censoring
assumption, support and model specification. This construction does not
contradict a small-interval approximation; it shows why that approximation
must be justified before applying the convention to coarse monthly data.
\citep{rel25pygwitness,rel25pygfit}

A competing event supplies different information. Consider 100 accounts with
20 prepayments, eight defaults and 72 accounts experiencing neither event.
With prepayment ordered first, its probability is $q=1/5$ and the conditional
default probability among the 80 nonprepaying accounts is $p=1/10$.
The factorization gives
\[
 \Pr(\text{default in the interval})=(1-q)p=\frac{2}{25}.
\]
This matches the selected outcome-fitting restriction and the simulator's
current competing-survival multiplier. Substituting the marginal default
frequency $2/25$ for the conditional hazard would apply that multiplier twice
and yield $8/125$. Consequently, the explicit competing-event branch is
coherent on this clock; it should not be repaired by indiscriminately putting
prepayment rows into the default denominator. The earlier two-period
comparison, $86/625$ against $19/100$ when prepayment is eliminated, retains
its distinct hypothetical targets. Exact records, event order and risk-set
denominators are the common basis for comparing Python and R implementations;
matching a method name or a plotted curve is insufficient.
\citep{rel25pygwitness,rel25pygfit}

